\subsection{Why the vacuity factors cannot simply be moved under the quantifier}
\label{sec:factor-movement-countermodel}

There is a concrete obstruction to replacing the original gate
\[
 F=\neg\bigl(\Delta\to(V\to\forall_c(S\to\neg P))\bigr)
\]
by
\[
 F_m=\neg\forall_c\bigl(\Delta\to(V\to(S\to\neg P))\bigr)
\]
using only the retained common heads.

\begin{lemma}[Failure of common-core factor movement]
Neither implication $F\to F_m$ nor $F_m\to F$ is uniformly
derivable from the retained common heads, even under $c\ne a,b$.
This does not decide interderivability of the complete gate theories.
\end{lemma}
\begin{proof}

Use worlds $(\eta,e)$, where $\eta:\mathbb N\to\{0,1\}$ and
$e\in\{0,1\}$. Interpret implication classically and define
\begin{align*}
 E(i,j)&\quad\text{iff}\quad\eta(i)=\eta(j),\\
 M(i,j)&\quad\text{iff}\quad
     e=0,\ \eta(i)=0,\ \eta(j)=1,\\
 \forall_i X&\quad\text{iff}\quad
     X\text{ holds at both }(\eta[i:=v],1-e),\quad v\in\{0,1\}.
\end{align*}
Evaluate natural-number formula codes by their exact PA196
constructors, with code zero true. Since $M(0,0)$ is always false,
code one is a genuine falsity in this interpretation. The alias
$I(0,0)=0$ respects classical implication.

Each quantifier accessibility relation is symmetric: a permitted
change of coordinate $i$ and color can be reversed by restoring that
coordinate and flipping the color again. Hence h1,
$\exists_i\forall_i X\to X$, is valid. Two different quantifier
relations commute because their coordinate reassignments commute
and two color flips cancel; this validates h2. Equality is actual
equality at a world, so the two membership-congruence heads h3 and
h4 hold. To validate h5, choose the new value of coordinate $i$ to
equal the current value of its witness coordinate; when the two
indices coincide the equality already holds. Head h6 is ordinary
normal-quantifier distribution. The restricted implicational head
is a classical tautology. MP and generalization preserve validity.
Thus every instance of every retained common head is valid.

Now take numerical parameters $a=0$, $b=2$, $c=1$, and an environment
with $\eta(0)=0$, $\eta(2)=1$. Both guard inequalities hold. Since
the property code $c=1$ is false, $P=(c\to M(c,b))$ is true at every
world. Also $\Delta$ is true. The membership-vacuity factor $V$ is
false at color zero and true at color one. The separation condition
$S=\forall_b((M(a,b)\to P)\wedge(M(c,b)\to c))$ is true at color
zero, and at color one it is true exactly when $\eta(c)=1$.
Consequently
\[
\begin{array}{c|cc}
 & e=0 & e=1\\\hline
 F & \text{false} & \text{true}\\
 F_m & \text{true} & \text{false}
\end{array}
\]
The two implications therefore both fail in a model of the common
core.
\end{proof}

The accompanying \file{core_countermodel.py} checks the binary-domain,
three-label instance exactly. Its 16 worlds suffice for the displayed
counterinstance; it checks h1 and h2 for all 65,536 truth sets, all
three-label instances of the atomic heads, and all local truth
assignments for h6. The preceding argument extends the retained-head
validity to arbitrary natural-number labels. A different gate that
moves these factors would need a new bootstrap argument establishing
the necessary vacuity properties; the successful conjunction-
distribution rewrite needs no such step.
